\documentclass[11pt,a4paper]{article}
\usepackage[margin=22mm]{geometry}
\usepackage{amsmath,amssymb,booktabs,microtype}
\usepackage[hidelinks]{hyperref}
\newcommand{\dd}{\mathrm d}
\newcommand{\tr}{\operatorname{tr}}
\newcommand{\ii}{\mathrm i}
\title{Route G-R7: One Clock Action,\\Frequency, Phase and Path Correlations}
\author{A conservative quantum extension of the G-R6 baseline}
\date{29 September 2026}
\begin{document}
\maketitle
\begin{abstract}
Preserve the E2/E3 energy gaps, thermal response and metric of G-R4--G-R6.
Promote the path label to a coherent two-branch degree of freedom.
The same proper-time action gives frequency, relative internal phase,
and ideal interferometric visibility. An explicit pure/mixed counterexample
shows why these three observations alone need not certify entanglement.
Joint observables distinguish the cases. The full scalar interferometer
phase, radiative environment and apparatus closure remain independent
missing inputs; no time law, experimental signal or hardware is invented.
\end{abstract}

\section{One conservative effective action}
The retained local Hamiltonian, in the basis $(g,e_2,e_3)$, is
\begin{equation}
 H_C(T)=\sum_{i=2,3}h\nu_i^0[1+b_i(T)]|e_i\rangle\langle e_i|,
 \qquad b_i(T)=\beta_i u(T),\quad u(T)=a_{\rm rad}T^4 .
\end{equation}
Here $\beta_i=-\Delta\alpha_i/(2h\epsilon_0\nu_i^0)$ and all pinned
G-R4 inputs are unchanged. The static scalar dipole approximation is
retained, not promoted to a complete radiative model.
For a prescribed branch $r=A,B$, use the effective action
\begin{equation}\label{eq:action}
 S_r=S_{{\rm guide},r}+\int_r
 \left[\ii\hbar\chi^\dagger\partial_\tau\chi
       -E_0(x)-\chi^\dagger H_C(T(x))\chi\right]\dd\tau .
\end{equation}
$E_0$ includes ground-state rest energy and any physical common
ground-level shift. G-R4's differential polarizabilities do not determine
that ground shift separately. Guide and pulse phases are explicit inputs,
not derived from the thermal gap data. We assume state-independent
guides and coincident final wave packets at the retained order; otherwise
the branch evolution must include the external motion as well.

Variation of the internal state gives $\ii\hbar\partial_\tau\chi=H_C\chi$.
For diagonal adiabatic evolution during the analyzed hold,
\begin{equation}
 U_r=\operatorname{diag}(1,e^{-\ii\theta_{2r}},e^{-\ii\theta_{3r}}),
 \quad \theta_{ir}=2\pi\nu_i^0\int_r[1+b_i(T_r)]\dd\tau_r .
 \label{eq:phase}
\end{equation}
Writing $\lambda_r=\dd\tau_r/\dd t$, the received frequency against
reference $R$ is $\nu_{i,r}^{(R)}=(\lambda_r/\lambda_R)\nu_i^0(1+b_i)$,
recovering G-R6. Thus frequency is the phase derivative, not a second
adjustable law. General paths require all segments, kinematic terms,
preparation phases and pulses. The fixed reference is not a third
entangled quantum system here. This extension uses the established
clock-interferometry framework~\cite{zych,roura}, not a new quantization rule.

\newpage
\section{Joint evolution fixes phase and visibility}
For a balanced path state and initial internal state $\rho_C$, the joint
evolution is controlled by the path. Define
\begin{equation}
 \mathcal U=|A\rangle\langle A|\otimes e^{\ii\zeta_A}U_A
            +|B\rangle\langle B|\otimes e^{\ii\zeta_B}U_B,
 \quad \zeta=\zeta_A-\zeta_B .
\end{equation}
$\zeta_r$ includes the scalar propagation, ground-energy, guide and
beam-splitter contributions. Tracing the clock gives
\begin{equation}\label{eq:reduced}
 \rho_P=\frac12
 \begin{pmatrix}1&e^{\ii\zeta}\kappa\\e^{-\ii\zeta}\kappa^*&1\end{pmatrix},
 \qquad \kappa=\tr(\rho_C U_B^\dagger U_A).
\end{equation}
Projection onto $(|A\rangle+e^{\ii\varphi}|B\rangle)/\sqrt2$ yields
\begin{equation}\label{eq:port}
 P_+(\varphi)=\tfrac12[1+\Re(e^{\ii(\varphi+\zeta)}\kappa)],
 \quad \mathcal V=|\kappa|,\quad
 \phi_{\rm fringe}=\zeta+\arg\kappa .
\end{equation}
The fringe phase is undefined at $\kappa=0$. Unequal path populations
$p_A,1-p_A$ instead give $\mathcal V=2\sqrt{p_A(1-p_A)}|\kappa|$.
Only the internal contribution $\arg\kappa$ is fixed without $\zeta$.
Moving a scalar term $s_r(\tau)I$ into $H_C$ and out of $E_0$ changes
$U_r$ and $\zeta_r$ oppositely; the product in Eq.~\eqref{eq:port}
is invariant. One cannot simply discard physical branch-dependent
ground energies as an unobservable overall phase.

\paragraph{Separate two-level preparations.}
Let $\delta_i=\theta_{iA}-\theta_{iB}$ and
$|\chi_0\rangle=\sqrt{1-p}|g\rangle+\sqrt p\,e^{\ii\eta}|e_i\rangle$.
Then
\begin{align}
 \kappa_i&=(1-p)+pe^{-\ii\delta_i},\\
 |\kappa_i|^2&=1-4p(1-p)\sin^2(\delta_i/2).
\end{align}
For $p=1/2$ this becomes
\begin{equation}
 \kappa_i=e^{-\ii\delta_i/2}\cos(\delta_i/2),\quad
 \mathcal V_i=|\cos(\delta_i/2)|,\quad
 1-\mathcal V_i=\delta_i^2/8+O(\delta_i^4).
\end{equation}
An energy eigenstate has visibility one: a different scalar phase is
not itself a which-path record in the internal state.
For a pure initial clock the joint state has Schmidt eigenvalues
$(1\pm|\kappa|)/2$, hence negativity
\begin{equation}
 \mathcal N=\frac{\sqrt{1-|\kappa|^2}}2 .
\end{equation}
This inference requires purity and the declared closed unitary model.

\paragraph{One shared ground is not two clock qubits.}
For a coherent equal three-level preparation,
\begin{equation}
 \kappa_{023}=\frac{1+e^{-\ii\delta_2}+e^{-\ii\delta_3}}3,
\end{equation}
not $\kappa_2\kappa_3$. E2/E3 runs or a genuine qutrit preparation must
be specified. No independent tensor factor is invented for each gap.

\newpage
\section{A counterexample and an actual entanglement discriminator}
Since $U_B^\dagger U_A$ is diagonal in this model,
\begin{equation}
 \kappa=\sum_j(\rho_C)_{jj}e^{-\ii\delta_j},\qquad\delta_g=0.
\end{equation}
Erasing initial energy coherences leaves the entire reduced path signal
unchanged, but can remove joint entanglement. A sharp two-level fixture is
$U_A=I$, $U_B=Z$, with
\begin{align}
 |\Psi_{\rm pure}\rangle
 &=\tfrac12(|00\rangle+|01\rangle+|10\rangle-|11\rangle),\\
 \rho_{\rm mixed}
 &=\tfrac12|+\rangle\langle+|\otimes|0\rangle\langle0|
   +\tfrac12|-\rangle\langle-|\otimes|1\rangle\langle1| .
\end{align}
The inputs are respectively $|+\rangle$ and $I/2$ for the clock.
Both output path states equal $I/2$, so both have zero visibility.
The pure state has negativity $1/2$; the explicitly separable mixture
has negativity zero. Frequencies and the diagonal conditional phase
map are the same. Visibility alone therefore does not certify entanglement.

One joint witness for this calibrated fixture is
\begin{equation}
 W=\tfrac12(I-X_P\otimes Z_C-Z_P\otimes X_C).
\end{equation}
For product Bloch vectors $a,b$,
$a_xb_z+a_zb_x\leq\sqrt{a_x^2+a_z^2}\sqrt{b_z^2+b_x^2}\leq1$;
convexity extends $\langle W\rangle\geq0$ to separable states.
The pure fixture gives $-1/2$ and the mixed fixture gives zero.
This is a joint-measurement specification, not a realized $\pi$-phase
experiment or a witness automatically optimal for a small phase.

\paragraph{State-resolved fringes and their limitation.}
With populations $p_j$ unchanged in both arms, joint port/energy
probabilities are
\begin{equation}
 P(+,j)=\frac{p_j}{2}
       [1+\cos(\varphi+\zeta-\delta_j)].
\end{equation}
Differences between excited and ground fringe phases isolate $-\delta_i$
if external phases are state-independent and pulses calibrated. This is
linear phase information, but still does not distinguish the pure/mixed
counterexample. Clock-coherence measurements and joint noncommuting
observables, or tomography, are needed for an entanglement claim.

\paragraph{Reversibility is not an arbitrary contrast correction.}
A common final clock unitary $V$ preserves $\kappa$ because
$(VU_B)^\dagger VU_A=U_B^\dagger U_A$.
More generally a local trace-preserving clock channel cannot change
$\rho_P$, by the defining property of the partial trace. Dephasing the
clock after the pure fixture makes it separable without changing the
path signal. Thus a clock's Ramsey decay cannot automatically be
multiplied into path visibility. A path-controlled inverse can undo the
ideal correlations; implementing it requires controls. Conditional
eraser signals must retain their selection probabilities.

\newpage
\section{Numerical hold kernel and a cross-transition consistency relation}
Use G-R6's illustrative 1 m height contrast and 300/310 K baths.
Let $t_H=10$ ms be a constant-setting hold, in reference proper-time
units, with $\lambda_B=1$, $\lambda_A=1+\gamma$ and
$\gamma=g_*\Delta h/c^2$. This is a prescribed hold contribution, not a
complete closed interferometer. To avoid subtracting nearly equal
optical frequencies, evaluate the algebraically equivalent expression
\begin{equation}
 \Delta\nu_i=\nu_i^0[\gamma+(b_{iA}-b_{iB})+\gamma b_{iA}],
 \qquad \delta_i=2\pi\Delta\nu_i t_H.
\end{equation}
The $\gamma b_{iA}$ product is retained as ordinary multiplication;
the weak-field background itself is not an all-orders metric.
\begin{center}
\begin{tabular}{llrrr}
\toprule
Hold contrast & Clock & $\Delta\nu$ (Hz) & $\delta$ (rad) & $1-\mathcal V$\\
\midrule
Gravity only & E2 & $0.07511$ & $0.004719$ & $2.784\,10^{-6}$\\
 & E3 & $0.07006$ & $0.004402$ & $2.422\,10^{-6}$\\
Thermal only & E2 & $-0.05051$ & $-0.003173$ & $1.259\,10^{-6}$\\
 & E3 & $-0.006500$ & $-0.0004084$ & $2.085\,10^{-8}$\\
\bottomrule
\end{tabular}
\end{center}
These are ideal two-level-preparation predictions, not observations
or achieved precision. The executable also includes identical paths,
both controls together, and a genuine shared-ground qutrit.
Frequency and phase retain the effect's sign; visibility alone is even
in $\delta$ and cannot determine that sign.

\paragraph{The same phase data must obey a thermal-rejection identity.}
Within the retained static response, define
\begin{equation}
 x_i=\frac{\delta_i}{2\pi\nu_i^0},\qquad
 J=\int[\lambda_Au_A-\lambda_Bu_B]\dd t .
\end{equation}
Equation~\eqref{eq:phase} gives, without a new slowdown function,
\begin{equation}
 x_i=\Delta\tau+\beta_iJ,\qquad
 \Delta\tau=\int(\lambda_A-\lambda_B)\dd t.
\end{equation}
The same G-R5 weights $w_2=\beta_3/(\beta_3-\beta_2)$ and
$w_3=-\beta_2/(\beta_3-\beta_2)$ therefore give
\begin{equation}\label{eq:phasecombination}
 w_2x_2+w_3x_3=\Delta\tau,\qquad
 x_3-x_2=(\beta_3-\beta_2)J.
\end{equation}
This is a combination of calibrated, unwrapped relative phases, not
visibilities. It needs matched trajectories and pulse conventions in
the E2/E3 preparations, independent radiometry, and propagated
polarizability uncertainty. Two channels cannot independently fit two
unknowns and simultaneously validate the law. At matched 300 K and
the illustrative hold, the reconstructed $\Delta\tau$ is
$1.0911\,10^{-18}$ s.

The frequency, phase and visibility predictions are distinct
observations constrained by one evolution operator, not independent
fit functions. In particular a thermally produced gap shift can also
produce path-clock correlation. Its mere observation would not
uniquely identify gravitational proper time or quantum gravity.

\newpage
\section{What remains physical input, and the next useful test}
\paragraph{A hold is not automatically a closed interferometer.}
The full paths, guiding potentials, launch/recombination dynamics,
initialization surfaces and laser phases must be specified.
State-dependent guiding and final packet mismatch change the simple
two-branch reduction. A standard freely falling light-pulse geometry
can cancel the putative uniform-field redshift sensitivity; one cannot
substitute $g\Delta h\,t/c^2$ without checking the whole sequence.
Adjustable clock initialization and differential phase protocols offer
established alternatives~\cite{roura}; their neutral-atom pulse design
is not assumed to work unchanged for this charged-ion E2/E3 system.
The prior G-R6 comparison of two sites is not preparation of one
ion in a coherent spatial superposition.

\paragraph{The conservative shift does not determine an open-system visibility.}
For full branch unitaries $W_A,W_B$ on clock and environment,
the cross amplitude is
\begin{equation}
 K=\tr_{C,E}(W_A\rho_{CE}W_B^\dagger).
\end{equation}
Only with initially factorized states and branch factorizations
$W_r=U_r\otimes V_r$ does this reduce to $\kappa$ times a separate
environment overlap. Photon which-path records, recoil, thermal
correlations and emission generally require the fuller model.
Real static polarizabilities constrain a conservative dispersive
shift; they do not specify the necessary dissipative correlation
functions. No fitted dephasing factor or decay rate is inserted here.

\paragraph{Do not extend the E2 idealization to seconds.}
The gravity-only first ideal zeros at 1 m would require
$t_{\pi,i}=[2\nu_i^0\gamma(1+b_i)]^{-1}$, approximately
6.66 s for E2 and 7.14 s for E3. Yet the E2 upper state has a
tens-of-milliseconds lifetime: a same-isotope study reports
61.8 ms with substantial systematic variation~\cite{schacht}.
A newer $^{174}$Yb$^+$ study~\cite{shao} concerns another isotope.
We do not pool these values or assign a precise $^{171}$Yb$^+$
environmental rate. Even 10 ms is only a conservative hold illustration,
not a decay-free experimental claim. E3's long-lived level would not
by itself establish coherent spatial splitting, guide closure or
laser coherence.

\paragraph{Bounded conclusion and priority.}
The unified conservative kernel, pure/mixed counterexample and
joint witness are complete under their stated assumptions, with
executable density-matrix checks. The missing physical controls are
recorded in \texttt{CLOCK\_PATH\_CARD.json}.
Next specify a closed path and calibrated state-resolved readout,
prioritizing the linear differential phase. Use coherent/dephased
preparations and joint measurements to substantiate entanglement.
The small-phase visibility loss is quadratic, so it is a consistency
check rather than a uniquely diagnostic discovery channel. No data,
lab contact, quantum-gravity detection, Ed law or Route-F promotion
is supplied by this mathematical extension.

\begin{thebibliography}{9}\small
\bibitem{zych} M. Zych, F. Costa, I. Pikovski and C. Brukner,
\href{https://doi.org/10.1038/ncomms1498}{Nature Communications 2, 505 (2011)}.
\bibitem{roura} A. Roura,
\href{https://doi.org/10.1103/PhysRevX.10.021014}{Physical Review X 10, 021014 (2020)}.
\bibitem{schacht} M. Schacht and M. Schauer,
\href{https://arxiv.org/abs/1310.2530}{arXiv:1310.2530 (2013)}, same-isotope lifetime study.
\bibitem{shao} H. Shao et al.,
\href{https://doi.org/10.1103/PhysRevResearch.5.023193}{Physical Review Research 5, 023193 (2023)}.
\end{thebibliography}
\end{document}
